\section{Reprezentacja stanu}
\label{sec:reprezentacja-stanu}

W badaniach zastosowano dwa warianty reprezentacji stanu.
Oba powstają wyłącznie na podstawie obiektu
\texttt{UTHObservation}, a więc informacji dostępnych graczowi.
Karty krupiera oraz pozostałe elementy pełnego stanu środowiska
nie są przekazywane do modeli.

Pierwszy wariant, dalej określany jako State A, stanowi reprezentację
surową. Dwie karty prywatne oraz pięć miejsc przeznaczonych na karty
wspólne kodowane są one-hot w przestrzeni 52 kart. Nieodkryte jeszcze
karty wspólne reprezentowane są wektorami zerowymi. Do reprezentacji
dołączono kod aktualnej fazy rozdania oraz maskę sześciu możliwych
akcji. Łączny wymiar State A wynosi

\begin{equation}
    d_A = 7 \cdot 52 + 3 + 6 = 373.
\end{equation}

State A nie zawiera ręcznie przygotowanych informacji o sile układu
pokerowego. Model musi więc samodzielnie nauczyć się zależności
pomiędzy widocznymi kartami a wartością poszczególnych decyzji.

Drugi wariant, State B, rozszerza State A o 20 cech domenowych.
Obejmują one właściwości kart prywatnych, kategorię najlepszego
widocznego układu oraz informacje opisujące strukturę kart,
między innymi pary, zgodność kolorów i możliwość utworzenia strita.
Wszystkie cechy wyznaczane są wyłącznie na podstawie kart widocznych
dla gracza.

Wymiar drugiego wariantu wynosi

\begin{equation}
    d_B = 373 + 20 = 393.
\end{equation}

Porównanie obu reprezentacji pozwala ocenić, czy przekazanie modelowi
jawnych cech pokerowych ułatwia proces uczenia w porównaniu
z reprezentacją opartą głównie na surowym kodowaniu kart.